% !PRINTED none

\begin{center}
\gesperrt{Sechster Abschnitt.}

\gesperrt{Tschirnhausen-Transformation.}

\rule{4cm}{0.4pt}

\S.\,68.

\gesperrt{Die Hermite'sche Form der Tschirnhausen-Transformation.}
\end{center}

Wir haben im vierten Abschnitt den Grundgedanken der
Tschirnhausen-Transformation schon kennen gelernt.

Die Aufgabe war die, eine algebraische Gleichung $n^{\text{ten}}$
Grades
\weq{1}{f(x) = a_0 x^n + a_1 x^{n-1} + \cdots + a_{n-1} x + a_n = 0}
durch eine Substitution $(n-1)^{\text{ten}}$ Grades
\weq{2}{y = \alpha_0 + \alpha_1 x + \alpha_2 x^2 + \cdots + \alpha_{n-1} x^{n-1}}
umzuformen, um in den willk\"urlichen Co\"efficienten $\alpha$ Mittel zu
gewinnen, die Gleichung zu vereinfachen.

Hermite hat dadurch, dass er der Substitution (2) eine
besondere Form gab, diese Aufgabe sehr vereinfacht und mit
der Invariantentheorie in Verbindung gebracht\footnote{\gesperrt{Hermite}, Sur quelques th\'eor\`emes d'alg\`ebre et la r\'esolution de
l'\'equation du quatri\`eme degr\'e, aus den Comptes rendus der Pariser Akademie besonders erschienen. Paris 1859.}.

Wir haben in \S.\,4 eine Reihe von Functionen kennen gelernt,
$f_0, f_1, f_2 \ldots f_{n-1}$, durch die sich die Potenzen von $x$ bis zur
$(n-1)^{\text{ten}}$ rational ausdr\"ucken lassen, und diese Functionen sind
es, die Hermite zur Darstellung der Substitution (2) verwendet.

Wir bezeichnen hier mit $x$ eine Wurzel der Gleichung (1),
w\"ahrend unter $t$ eine unbestimmte Ver\"anderliche verstanden sein
soll. Dann ist nach \S.\,4
%CONT>>
